\section*{ID7089: How This Compares to Standard Theory: α}
\paragraph{Metadata.}
PiID: 3904688383436 | RTEID: RTE:P27528.1613:T1.3756 | UUID: 216ec34a-bb71-559d-93db-1d731f692cde

\paragraph{Canonical Statement.}
\# 7) How this compares to standard theory - **Newtonian gravity / GR:** predict no dependence of gravitational acceleration on the test mass (universality of free fall). That corresponds to \textbackslash{}(\textbackslash{}alpha=0\textbackslash{}). - **If you measure a nonzero α (significantly outside uncertainties)** you would have discovered a violation of equivalence principle / new coupling. The standard experimental literature places very tight constraints on composition-dependent violations (Eötvös parameter limits \textbackslash{}( \textbackslash{}sim 10\textasciicircum{}\{-13\}\textbackslash{}) or better), but your α parameter as defined here is suppressed by \textbackslash{}(m/M\_\textbackslash{}oplus\textbackslash{}) which makes an α of moderate size produce tiny Δa unless α is enormous.

\paragraph{Canonical Equation.}
\[
\alpha=0
\]

\paragraph{Dictionary.}
Relation: α=0

\paragraph{Encyclopedia.}
How This Compares to Standard Theory: α is an algebraic definition, written α=0. It is an algebraic definition expressing the quantity directly in terms of the variables on its right-hand side. Within the Trawin Zero Point registry it serves as one element of the framework's mathematical narrative.
